\begin{resumenApertura}
audIT oferta la implementación y la operación de la plataforma de Transportes Curimón como un
contrato de alcance completo, con precios por fase y pagos contra aceptación. Los valores se
concilian con la planilla de presupuesto 3.2 y sus partidas, sin incorporar recuperos hipotéticos
del mandante a los cobros de audIT.
\begin{recibe}
\item Precio, hitos, servicios incluidos, exclusiones y condiciones comerciales de la oferta.
\end{recibe}
\end{resumenApertura}

\section{Propuesta de valor económico}\label{sec:eco-valor}

La \tabref{eco-valores} presenta los precios de implementación y de los 36 meses de operación.
La suma es el valor total del contrato. Los valores netos, IVA y total se muestran en las tres
monedas del \form{E-24}{73}; la fase de operación no se anticipa a implementación.

\tablaMontos{Valores ofertados por fase}{eco-valores}{implementacion,operacion}

El precio total neto es \montoPartida{proyecto}. El hardware de campo y de red adquirido por el
mandante se informa por separado: no incrementa el ingreso ni el egreso de audIT.

\subsection{Tipo de cambio, reajustabilidad y vigencia}\label{sec:eco-condiciones}

Para esta oferta de fecha 25 de noviembre de 2026 se utilizan los parámetros convencionales
del E-24. Esa fecha identifica la aplicación de las bases a la oferta, no una cotización de
mercado. Los precios de ambas etapas de implementación son firmes, conforme al
\art{9.4}{7}. La operación se oferta sin reajuste: audIT asume la variación de sus costos
dentro del precio comprometido. El presupuesto base no incorpora reajustes de ingresos.

La oferta tiene vigencia comercial de 150 días corridos contados desde la fecha de entrega.
La vigencia y condiciones de la garantía de seriedad se acreditan por separado en el Sobre
N.º 1; este plazo comercial no sustituye dicha garantía ni certifica su emisión.

\section{Estructura de pagos detallada}\label{sec:eco-pagos}

La implementación se factura contra los doce hitos del \form{E-25}{74 y 75}. Un hito
presupuestado no acredita un acta ni un pago realizado: la aceptación exige el pronunciamiento
de la Contraparte Técnica y las evidencias del \art{18}{13}.

\subsection{Hitos de implementación}\label{sec:eco-hitos}

La \tabref{hitos-e25} distribuye el precio de implementación. Los porcentajes son los de
las bases y suman el 100 por ciento; los importes provienen de la partida implementación.

\tablaHitosPago

La \tabref{eco-criterios} vincula cada hito con la evidencia que debe aceptar la contraparte.
Los protocolos detallados y sus trazas a requerimientos se conservan en T-13 y T-17.

\begin{tabla}[ancho=completo]{Aceptación vinculada a los hitos de pago}{eco-criterios}{L{12mm} X}
Hito & Evidencia y condición de aceptación \\ \midrule
H1 & Acta de alcance, catálogo y matriz de trazabilidad aprobados. \\
H2 & Arquitectura, seguridad y modelo de datos aprobados. \\
H3 & Infraestructura híbrida y cuatro ambientes habilitados, con observabilidad probada. \\
H4 & Candidato E1 con QA superado y evidencia de cobertura, listo para pruebas. \\
H5 & Aceptación de usuario, carga, resiliencia y seguridad ofensiva E1 aprobadas. \\
H6 & Inicio de marcha blanca E1, reversión operativa y usuarios capacitados. \\
H7 & Cierre conforme de marcha blanca E1 y acta de paso a producción. \\
H8 & Línea base y diseño detallado E2 aprobados. \\
H9 & Candidato E2 con QA superado, listo para pruebas. \\
H10 & Certificación E2 y cierre del desarrollo aceptados. \\
H11 & Inicio de marcha blanca E2 en convivencia con E1 en producción. \\
H12 & Paso a producción E2 y aceptación final de implementación. \\
\end{tabla}

Cada acta debe acompañarse de artefacto, verificación objetiva, trazabilidad y resolución de
observaciones. Se respetan los plazos de revisión y subsanación del artículo 18; el modelo
distingue fecha contractual, facturación y cobro efectivo.

\subsection{Pagos mensuales de operación}\label{sec:eco-mensual}

La operación se factura de M21 a M56 en 36 pagos mensuales. El caso base cobra cada servicio
en el mes siguiente, de M22 a M57. La mensualidad neta es \montoPartida{cobro-m57}.
El componente fijo incluye operación, soporte, mantención, nube, conectividad y los servicios
presupuestados. No se presupuesta un componente variable positivo adicional ni cobros por hora
fuera de ese precio. Las deducciones por incumplimiento de SLA se aplican conforme al
\art{78}{40} y al E-25. La revisión de la factura y los indicadores no sustituye el
procedimiento de aplicación de multas.

\section{Resumen ejecutivo de valor}\label{sec:eco-resumen}

La \tabref{eco-incluidos} delimita los productos y servicios comprendidos. La asignación de
pagadores se concilia con T-11 y las hojas de proveedores para evitar que una exclusión
comercial elimine una prestación ya costeada.

\begin{tabla}[ancho=completo]{Servicios y productos incluidos}{eco-incluidos}{L{35mm} X}
Componente & Inclusión contractual \\ \midrule
Plataforma & Infraestructura on-premise de audIT, nube, observabilidad, seguridad y ambientes. \\
Software y datos & Alcances E1/E2, integraciones, migración, pruebas y documentación. \\
Instalación & Piloto, montaje externo por lotes M7--M9, banco, integración, supervisión y recepción interna. \\
Adopción & Capacitación, acompañamiento de marchas blancas y transferencia. \\
Operación & Mesa de ayuda, guardias, mantenimiento y campo externo por visita con coordinación interna. \\
Innovaciones & Cinco propuestas S13/T-19, con costo y atribución de personal identificados. \\
\end{tabla}

La infraestructura, integración y atención permanecen bajo responsabilidad de audIT aunque se
contraten servicios externos. Los criterios de cobertura y recepción son los de S3, S7 y S9.

\subsection{Exclusiones}\label{sec:eco-exclusiones}

Curimón adquiere el hardware de campo y de red especificado en T-11. Su referencia presupuestaria
es \montoPartida{hardware-mandante}, informada fuera del flujo de audIT. La exclusión de
combustible corresponde a la flota vehicular: no excluye las pruebas del grupo electrógeno
presupuestadas por audIT. No se eliminan de la oferta la energía, infraestructura, instalación,
conectividad ni servicios que el modelo incluye a cargo de audIT. Las ampliaciones de alcance
requieren gestión formal de cambios.

\subsection{Dependencias comerciales}\label{sec:eco-supuestos}

La puesta en servicio requiere equipos homologados y recibidos, accesos a terminales, ventanas
de montaje y factibilidad de integraciones con terceros. Las tarifas y calendarios de proveedores
son supuestos presupuestarios hasta contar con respaldo comercial. El financiamiento inicial y
la facilidad contingente no se presentan como aprobaciones bancarias emitidas. La subcontratación
se ejecutará bajo las autorizaciones y límites del \art{73}{38}, manteniendo responsabilidad
de audIT sobre el servicio.

\subsection{Beneficios del mandante}\label{sec:eco-beneficios}

El beneficio esperado se mide contra líneas base operacionales de S13/T-19. El recupero bruto
referencial de INN-4 es\par\noindent\montoPartida{recupero-inn4} en 36 meses, bajo el supuesto de recuperar
la mitad del monto objetado. No equivale a beneficio neto ni a cobro de audIT. Las otras mejoras
se someten a medición incremental y adopción, sin incorporar ahorros hipotéticos a ingresos
contractuales ni duplicar costos compartidos.

\subsection{Términos comerciales}\label{sec:eco-terminos}

La oferta comprende el alcance indivisible y el calendario del \art{17}{12}. Se expresan
neto, IVA y total en las monedas de evaluación. La facturación sigue el E-25 y la aceptación
del artículo 18. Las ampliaciones, exclusiones o cambios deben formalizarse sin alterar de
manera unilateral los plazos ni los criterios de recepción. Las referencias a financiación,
proveedores y beneficios mantienen sus condiciones de verificación y no acreditan ejecución.
